\item Considere aire (un gas ideal diatómico) a $27.0^\circ\text{C}$ y presión atmospheric que entra a una bomba de bicicleta que tiene un cilindro con diámetro interno de 2.50 cm y longitud 50.0 cm. La carrera hacia abajo comprime adiabáticamente el aire que entra a la llanta. Se desea investigar el incremento de temperatura de la bomba. 
(a) ¿Cuál es el volumen inicial de aire en la bomba? 
(b) ¿Cuál es el número de moles de aire en la bomba? 
(c) ¿Cuál es la presión absoluta del aire comprimido? 
(d) ¿Cuál es el volumen del aire comprimido? 
(e) ¿Cuál es la temperatura del aire comprimido? 
(f) ¿Cuál es el incremento de energía interna del gas durante la compresión? 
\textbf{¿Qué pasaría si?} La bomba está hecha de acero con espesor de 2.00 mm. Suponga una longitud del cilindro de 4.00 cm en equilibrio térmico con el aire. 
(g) ¿Cuál es el volumen de acero en esta longitud de 4.00 cm? 
(h) ¿Cuál es la masa del acero en esta longitud de 4.00 cm? 
(i) Suponga que la bomba se comprime otra vez. Después de la expansión adiabática, la conducción resulta en el incremento de energía en el inciso (f) compartido entre el gas y la longitud de acero de 4.00 cm. ¿Cuál será el aumento en temperatura del acero después de una compresión?
